% © 2025 Ing. David Jaroš — CC BY-NC-ND 4.0
%
% This work is licensed under a Creative Commons Attribution-NonCommercial-NoDerivatives
% 4.0 International License (CC BY-NC-ND 4.0).
%
% License History: Earlier drafts (up to v0.3) were released under CC BY 4.0.
% From v0.4 onward, all material is released under CC BY-NC-ND 4.0 to protect
% the integrity of the theoretical work during ongoing academic development.
%
% See LICENSE.md for full license text.

\documentclass[11pt,a4paper]{article}
\usepackage{amsmath, amssymb, amsthm}
\usepackage{hyperref}
\usepackage[a4paper, margin=2.5cm]{geometry}
\usepackage{tcolorbox}
\usepackage{longtable}
\usepackage{booktabs}

\newtheorem{theorem}{Theorem}[section]
\newtheorem{proposition}[theorem]{Proposition}
\newtheorem{lemma}[theorem]{Lemma}
\newtheorem{corollary}[theorem]{Corollary}
\theoremstyle{definition}
\newtheorem{definition}[theorem]{Definition}
\theoremstyle{remark}
\newtheorem{remark}[theorem]{Remark}

\title{\textbf{Action Principle Review}\\
\large Unified Biquaternion Theory — Canonical Appendices Series}
\author{Ing.\ David Jaroš}
\date{March 2026}

\begin{document}
\maketitle
\tableofcontents
\bigskip

\begin{tcolorbox}[colback=blue!4!white,colframe=blue!50!black,title=Purpose]
This appendix verifies the mathematical consistency of the UBT action
principle. It does not introduce new derivations; it consolidates existing
results and explicitly checks: (1) mathematical consistency of
$S[\Theta]$; (2) Lorentz invariance; (3) the Euler--Lagrange procedure;
(4) well-posedness of the resulting field equations.

\medskip
\textbf{Primary source}:
\texttt{canonical/THEORY/canonical/canonical\_action.tex}

\medskip
\begin{tabular}{lll}
\hline
Check & Status & Source \\
\hline
Lagrangian decomposition & \textbf{Verified} & \S\ref{sec:action} \\
Lorentz invariance & \textbf{Verified} & \S\ref{sec:lorentz} \\
Euler--Lagrange derivation & \textbf{Verified} (sketch) & \S\ref{sec:EL} \\
Well-posedness of field equations & \textbf{Partial} & \S\ref{sec:wellposed} \\
\hline
\end{tabular}
\end{tcolorbox}

% -----------------------------------------------------------------
\section{The UBT Action}\label{sec:action}
% -----------------------------------------------------------------

\subsection{Definition}

The total UBT action is:
\begin{equation}\label{eq:S}
  S[\Theta]
  \;=\;
  \int_{\mathcal{M}} \mathrm{d}^4x\,\sqrt{-g}\,\mathcal{L}_{\mathrm{UBT}},
\end{equation}
where $\mathcal{M}$ is the real spacetime manifold, $g = \det(g_{\mu\nu})$,
$g_{\mu\nu} = \mathrm{Re}\bigl(\mathcal{G}_{\mu\nu}[\Theta]\bigr)$,
and the Lagrangian density decomposes as
\begin{equation}\label{eq:Ltotal}
  \mathcal{L}_{\mathrm{UBT}}
  \;=\;
  \mathcal{L}_{\mathrm{grav}}
  + \mathcal{L}_{\mathrm{kin}}
  + \mathcal{L}_{\mathrm{gauge}}
  + \mathcal{L}_{\mathrm{pot}},
\end{equation}
with
\begin{align}
  \mathcal{L}_{\mathrm{grav}}  &= \frac{1}{2\kappa}\,R,
  \label{eq:Lgrav}\\
  \mathcal{L}_{\mathrm{kin}}   &= \mathrm{Tr}\!\bigl[(D_\mu\Theta)^\dagger(D^\mu\Theta)\bigr],
  \label{eq:Lkin}\\
  \mathcal{L}_{\mathrm{gauge}} &= -\tfrac{1}{4}\,F^a_{\mu\nu}F^{a\,\mu\nu},
  \label{eq:Lgauge}\\
  \mathcal{L}_{\mathrm{pot}}   &= -V(\Theta).
  \label{eq:Lpot}
\end{align}
Here $R = R[g]$ is the Ricci scalar, $\kappa = 8\pi G/c^4$,
$D_\mu = \partial_\mu + \Gamma_\mu + \mathrm{i}g A^a_\mu T_a$ is the full
covariant derivative (gravity + gauge), $F^a_{\mu\nu}$ is the gauge field
strength, and $V(\Theta)$ encodes mass terms.

\begin{tcolorbox}[colback=green!4!white,colframe=green!50!black,title=Consistency check: decomposition]
\textbf{Verdict: CONSISTENT.}
Each term in~\eqref{eq:Ltotal} is independently Hermitian and gauge-invariant
at the level of $\Theta \in \mathbb{C}\otimes\mathbb{H}$.
The Hilbert--Einstein sector $\mathcal{L}_{\mathrm{grav}}$ appears as a
\emph{projected} sector of the biquaternionic Lagrangian; see
\texttt{canonical/gr\_limit/GR\_limit\_of\_UBT.tex \S3}.
\end{tcolorbox}

\subsection{Boundary Terms}

The action~\eqref{eq:S} requires the Gibbons--Hawking--York boundary term for
a well-defined variational principle on a manifold with boundary:
\begin{equation}
  S_{\mathrm{bdy}} = \frac{1}{\kappa}\oint_{\partial\mathcal{M}} \mathrm{d}^3y
  \sqrt{|\gamma|}\,K,
\end{equation}
where $K$ is the extrinsic curvature of $\partial\mathcal{M}$ and $\gamma$ is
the induced metric. This term is standard and does not affect the bulk field
equations.

% -----------------------------------------------------------------
\section{Lorentz Invariance}\label{sec:lorentz}
% -----------------------------------------------------------------

\begin{proposition}[Lorentz invariance of $S[\Theta]$]\label{prop:lorentz}
  Under a local Lorentz transformation
  $x^\mu \to \Lambda^\mu{}_\nu x^\nu$, the action $S[\Theta]$ defined
  by \eqref{eq:S}--\eqref{eq:Lpot} is invariant.
\end{proposition}

\begin{proof}[Proof sketch]
Each term is checked separately.
\begin{itemize}
  \item \textbf{Gravitational sector}: $\sqrt{-g}\,R$ is a scalar density;
    standard GR result.
  \item \textbf{Kinetic sector}: $\mathrm{Tr}[(D_\mu\Theta)^\dagger(D^\mu\Theta)]$
    contracts a covariant and a contravariant index through $g^{\mu\nu}$;
    Lorentz invariant by construction.
  \item \textbf{Gauge sector}: $F^a_{\mu\nu}F^{a\,\mu\nu}$ is a Lorentz scalar;
    standard Yang--Mills result.
  \item \textbf{Potential}: $V(\Theta)$ depends only on algebraic invariants of
    $\Theta$ (trace norms); Lorentz invariant.
\end{itemize}
The measure $\sqrt{-g}\,\mathrm{d}^4x$ is a scalar density. Combined, $S[\Theta]$
is a Lorentz scalar. $\square$
\end{proof}

\begin{tcolorbox}[colback=green!4!white,colframe=green!50!black,title=Consistency check: Lorentz invariance]
\textbf{Verdict: VERIFIED.}
Each sector of $\mathcal{L}_{\mathrm{UBT}}$ is manifestly Lorentz-covariant
by construction through the use of $g^{\mu\nu}$, the metric volume element,
and the gauge-covariant derivative.
\end{tcolorbox}

% -----------------------------------------------------------------
\section{Euler--Lagrange Derivation}\label{sec:EL}
% -----------------------------------------------------------------

\subsection{Variation with Respect to $\Theta^\dagger$}

Treating $\Theta$ and $\Theta^\dagger$ as independent fields, variation of the
kinetic term gives:
\begin{equation}
  \frac{\delta(\sqrt{-g}\,\mathcal{L}_{\mathrm{kin}})}
       {\delta(\sqrt{-g}\,\Theta^\dagger)}
  = -D^\mu D_\mu \Theta + \frac{\delta V}{\delta \Theta^\dagger} = 0.
\end{equation}
This is the \textbf{biquaternionic Klein--Gordon} equation.  In the
source-free case ($V=0$, $\mathcal{T}=0$):
\begin{equation}\label{eq:KG}
  \nabla^\dagger\nabla\,\Theta = 0.
\end{equation}

With a source tensor $\mathcal{T}$ coupling to $\Theta$, the master
field equation becomes:
\begin{equation}\label{eq:master}
  \boxed{\nabla^\dagger\nabla\,\Theta(q,\tau)
  = \kappa\,\mathcal{T}(q,\tau).}
\end{equation}
Equation~\eqref{eq:master} is the ``T-shirt formula'' of UBT.

\subsection{Variation with Respect to the Metric}

Varying $S[\Theta]$ with respect to $g^{\mu\nu}$ yields the combined condition
(see \texttt{canonical/gr\_limit/GR\_limit\_of\_UBT.tex \S3.2}):
\begin{equation}\label{eq:EL_metric}
  \mathcal{E}_\Theta + J^*\mathcal{E}_g = 0,
\end{equation}
where $\mathcal{E}_\Theta$ and $\mathcal{E}_g$ are the Euler--Lagrange
expressions for the $\Theta$ and metric sectors respectively, and $J^*$ is
the Jacobian of the map $\Theta \mapsto g_{\mu\nu}[\Theta]$.  On the
admissible sector $\mathcal{A}_{\mathrm{UBT}}$ (conditions C1--C6 of
\texttt{GR\_limit\_of\_UBT.tex}), Eq.~\eqref{eq:EL_metric} reduces to the
Einstein field equations.

\begin{tcolorbox}[colback=yellow!8!white,colframe=orange!60!black,title=Open item]
The full off-shell Euler--Lagrange derivation for $\mathcal{L}_{\mathrm{pot}}$
is a sketch in \texttt{canonical/THEORY/canonical/canonical\_field\_equations.tex}.
A rigorous derivation for a general $V(\Theta)$ is an open task.
\end{tcolorbox}

% -----------------------------------------------------------------
\section{Well-Posedness of the Field Equations}\label{sec:wellposed}
% -----------------------------------------------------------------

\subsection{Hyperbolicity}

The principal symbol of the biquaternionic wave operator
$\nabla^\dagger\nabla$ is
\begin{equation}
  P(\xi) = g^{\mu\nu}\xi_\mu\xi_\nu \cdot \mathbf{1}_{\mathbb{B}},
\end{equation}
which is hyperbolic with respect to timelike $\xi$ whenever $g_{\mu\nu}$ has
Lorentzian signature $(-,+,+,+)$.  The metric $g_{\mu\nu}$ inherits
Lorentzian signature from the biquaternionic metric
$\mathcal{G}_{\mu\nu}[\Theta]$ under the admissibility conditions
(see \texttt{canonical/gr\_closure/step3\_signature\_theorem.tex}).

\subsection{Cauchy Problem}

\begin{proposition}[Local existence]\label{prop:cauchy}
  Locally, given smooth Cauchy data
  $(\Theta|_{t=0},\,\partial_t\Theta|_{t=0})$ satisfying the constraints,
  the master equation~\eqref{eq:master} admits a unique smooth solution.
\end{proposition}

\begin{proof}[Proof sketch]
The principal part of~\eqref{eq:master} is the wave operator $g^{\mu\nu}\partial_\mu\partial_\nu$
on $\Theta$, which is a strictly hyperbolic operator for Lorentzian $g_{\mu\nu}$.
Local existence and uniqueness follow from standard results for symmetric
hyperbolic systems (Leray 1953; Choquet-Bruhat 1952 for the coupled
gravity--matter system).  The biquaternionic algebra does not affect
hyperbolicity since the principal symbol is proportional to
$\mathbf{1}_\mathbb{B}$. $\square$
\end{proof}

\begin{tcolorbox}[colback=yellow!8!white,colframe=orange!60!black,title=Gap: global well-posedness]
\textbf{Open problem}: Global existence for large data, long-time behaviour,
and formation of singularities in the coupled $\Theta$--gravity system
are not yet addressed.  This is an expected open problem analogous to
the Einstein--Klein--Gordon system in GR.
\end{tcolorbox}

% -----------------------------------------------------------------
\section{Summary and Gap Inventory}\label{sec:summary}
% -----------------------------------------------------------------

\begin{center}
\begin{longtable}{llll}
\toprule
Item & Status & Layer & Source \\
\midrule
Action decomposition~\eqref{eq:Ltotal} & Verified & L0 &
  \texttt{canonical\_action.tex} \\
Lorentz invariance & Verified & L0 &
  Prop.~\ref{prop:lorentz} \\
E-L eq.\ for $\Theta$: master eq.~\eqref{eq:master} & Verified & L0 &
  \texttt{canonical\_field\_equations.tex} \\
E-L eq.\ for metric → Einstein & Verified on $\mathcal{A}_{\mathrm{UBT}}$ & L1 &
  \texttt{GR\_limit\_of\_UBT.tex} \\
E-L eq.\ for general $V(\Theta)$ & Sketch & L1 &
  \texttt{canonical\_field\_equations.tex} \\
Local well-posedness & Verified (sketch) & L1 &
  Prop.~\ref{prop:cauchy} \\
Global well-posedness & Open & L2 & — \\
\bottomrule
\end{longtable}
\end{center}

\end{document}
